% Adaptive attacks (C5): method descriptions, drafted 2026-09-05; prose revised 2026-09-05. Results are in attack_results.tex.
\section{Adaptive Attacks}\label{sec:attacks}

The audit so far used the risky model of the released logs, which cannot reproduce the passages at all. The adversary of Section~\ref{sec:threat} is allowed a stronger one, and Section~\ref{sec:memorising} builds it (checks C5 and C7). Sections~\ref{sec:composition} and~\ref{sec:attack-results} describe the attack and what it recovers, Section~\ref{sec:natural} repeats it with the risky model the mechanism's authors evaluate, Section~\ref{sec:cpfuse} applies the same attack to a mechanism with no budget, and Section~\ref{sec:prefixdebt} asks which part of the decoder does the protecting.

\subsection{A risky model that memorises}\label{sec:memorising}
Eq.~\eqref{eq:knaf} is stated for any risky model, and the adversary of Section~\ref{sec:threat} may choose it. Llama-3.1-8B-Instruct assigns the CopyBench passages $2.6$ nats per token (Section~\ref{sec:certificate}), so it cannot reproduce them and a decoder anchored to it has nothing to prevent. We use two risky models that can. The first is the adversary's: the 8B instruct model fine-tuned with low-rank adapters (rank 64 on every projection matrix, learning rate $2\times10^{-4}$, twelve epochs) on the 608 passages of the CopyBench attack and validation splits, presented as raw ``Complete the prefix'' text and in the chat template, with the 150 test passages held out; the adapters are merged, so the decoder sees an ordinary Llama-3.1-8B checkpoint. The recipe is released; the merged weights are not, because they reproduce copyrighted text. Section~\ref{sec:attack-results} verifies its memorisation without the anchor. The second is the deployer's: Llama~3.1~70B base, the risky model of the mechanism's own evaluation, which memorised two of the CopyBench books in pre-training~\cite{cooper2025extracting} (Section~\ref{sec:natural}).

\subsection{Composition: windowed extraction under compliant queries}\label{sec:composition}
Cohen's composition theorem~\cite{cohen2025blameless} says that a per-prompt near-access-freeness guarantee does not bound what an adversary learns from many prompts. We instantiate it against the deployed decoder. For each memorised passage the first $20$ tokens are the seed and the remaining tokens (mean $260$) the target, and three query strategies are compared at the same $k$. A \emph{single} query asks for the whole target with $T_{\max} = 296$, the longest target in the set, so every single query has budget $K = 296k$ with banking allowed. \emph{Oracle} windows request window $i$ of length $L \in \{20, 50\}$ with the true text up to its start as the prompt, the query pattern of Cohen's theorem and the position of a rights-holder testing their own work; \emph{chained} windows use the seed plus everything extracted so far, which needs no oracle. Every query is an ordinary call with budget $K = kL$, and one with $Z > \max\{0, B\}$ would disqualify the attack. Outputs are stitched and scored by near-verbatim recall and longest exact span, for $k \in \{0.15, 0.5, 1, 3, 5, 10, 20\}$ and both baselines. The prefix debt charged at the start of every query (Eq.~\eqref{eq:banking}) is the mechanism's only defence here, and its size on a memorised prefix is itself a measurement we report. A bank-and-burst variant could not be evaluated: neither memorising model follows a filler instruction, so nothing is banked and the attack degenerates to the single query. Section~\ref{sec:odometer} bounds bursts directly instead.
